\bookchapter{Capítulo 12: Entretenimento Audiovisual e Leitura}

\chapterintro{O celular hoje basicamente já substituiu o MP3, o controle remoto da TV, o Kindle e o console portátil. Mas, sinceramente, muita gente só usa as funções mais básicas desses apps, sem descobrir um monte de recursos úteis. Neste capítulo, vamos falar sobre como tornar o entretenimento no celular mais prático e mais confortável.}

\booksection{I. Ouvir Música: Qual App Escolher e Como Ouvir Confortavelmente}

Os apps de música hoje têm funções bem parecidas, mas cada um tem seu foco; escolher certo poupa muito trabalho:

\begin{itemize}
  \item \textbf{NetEase Cloud Music}: seu ponto forte é a comunidade de playlists e as seções de comentários. Se você quiser ouvir "músicas para dias de chuva" ou "música instrumental para relaxar depois de estudar", é só pesquisar que já tem playlists organizadas por outras pessoas. Os comentários às vezes são até mais interessantes que as músicas. A desvantagem é que parte do catálogo de músicas tem direitos autorais incompletos; algumas músicas ficam indisponíveis de repente.
  \item \textbf{QQ Music}: o catálogo de direitos autorais mais completo. Cantores populares como Jay Chou e Mayday basicamente só estão disponíveis no QQ Music. As opções de qualidade de áudio também são muito ricas; assinantes podem ativar a qualidade sem perdas. Se você se importa muito em "encontrar todas as músicas que quer ouvir", o QQ Music é a escolha mais segura.
  \item \textbf{Apple Music}: a primeira escolha para usuários Apple. Interface limpa, sem anúncios, boa qualidade de áudio, e ainda suporta áudio sem perdas e Dolby Atmos. Mas seu algoritmo de recomendação não entende tão bem as preferências do usuário chinês quanto os apps nacionais, e a cultura de playlists não é tão desenvolvida quanto a do NetEase. Ideal para quem gosta de experiência simples e não se importa muito com funções sociais.
  \item \textbf{Spotify}: o maior número de usuários no mundo, algoritmo de recomendação reconhecidamente o melhor. O "Discover Weekly" e o "Daily Mix" têm uma precisão assustadora; quanto mais você usa, mais ele te conhece. A desvantagem é que a versão gratuita tem muitos anúncios, e a cobertura de direitos autorais de algumas músicas chinesas não é completa. Requer rede especial para funcionar.
\end{itemize}

\textbf{Algumas dicas para ouvir música}: se você usa fone Bluetooth, veja nas configurações do app de música se há “Ajuste de Efeito Sonoro” ou “Equalizador”; mexer um pouco nos graves e agudos melhora bastante o resultado. Além disso, streaming de música consome muitos dados móveis; lembre-se de ativar nas configurações “Reproduzir sem perdas somente no Wi-Fi”, ou baixe antecipadamente suas playlists favoritas para ouvir offline.

\booksection{II. Podcasts: o “Companheiro de Ouvido” no Trajeto}

Podcast é uma coisa que quem nunca ouviu acha que não tem nada a ver, mas depois que começa a ouvir, não consegue mais parar. Na essência, é a “versão Netflix do rádio” — tem gente falando de tecnologia, de história, lendo histórias; há de todo tipo que você quiser ouvir.

\begin{itemize}
  \item \textbf{App Xiaoyuzhou}: o melhor app de podcast feito na China, interface limpa, sem funções bagunçadas. Sua "Recomendação Diária" e "Seleção do Editor" ajudam a descobrir novos programas, e a seção de comentários também é bem ativa. O que eu mais ouço pessoalmente é o "Huzuohuyou" (cultura e história) e "Buheshiyi" (temas sociais); cada episódio é como estar ouvindo pessoas inteligentes conversando.
  \item \textbf{Apple Podcasts}: nativo do iPhone, não precisa baixar app adicional. A melhor integração com o CarPlay, muito prático para dirigir.
  \item \textbf{Dicas de escuta}: a maioria dos apps de podcast suporta reprodução em 1,2x ou 1,5x; depois que você se acostuma, economiza bastante tempo. Eu geralmente ouço em 1,2x para chinês e 1,5x para inglês. Além disso, é essencial ativar "Baixar automaticamente novos episódios" e "Excluir automaticamente após ouvir"; assim, seu celular não fica cheio de arquivos de áudio já ouvidos.
\end{itemize}

\booksection{III. Assistir Vídeos: Do Player ao Streaming, o Jeito Certo de Usar}

Muita gente não sabe que o celular consegue reproduzir muito mais formatos do que os oferecidos nos apps de streaming. Um bom player local resolve muitos problemas.

\begin{itemize}
  \item \textbf{VLC Player}: open source e gratuito, reproduz praticamente qualquer formato — MKV, AVI, FLAC e até imagens de disco ISO. Se você baixou um vídeo que roda no computador mas não no celular, 80\% de chance de ser problema de formato; instale o VLC e resolve. Ele também suporta legendas externas; coloque o arquivo de legenda .srt na mesma pasta do vídeo que o VLC carrega automaticamente.
  \item \textbf{nPlayer}: se você quer uma experiência mais refinada, este app pago vale a pena. Interface bonita, controles por gesto muito práticos — deslizar para cima/baixo na esquerda ajusta o brilho, na direita ajusta o volume, deslizar para os lados avança/retrocede. Suporta DTS e som Dolby; ao assistir blockbusters, a sensação de imersão é muito maior.
  \item \textbf{Bilibili (site B)}: já não é só um site de cultura otaku; hoje a quantidade de recursos de aprendizado é absurda. Pós-graduação, programação, fotografia, culinária — tem tutorial de tudo. A cultura de danmaku (comentários flutuantes) é a alma do Bilibili; assistir com os danmaku ligados dá a sensação de que tem um monte de gente assistindo junto com você. Mas, ao ver vídeos de estudo, sugiro desligar os danmaku, senão é fácil se distrair.
  \item \textbf{Transmissão para a TV}: por maior que seja a tela do celular, não é tão confortável quanto a TV. No mesmo Wi-Fi, abra o botão "TV" no canto superior direito do app de vídeo e projete a imagem na televisão. A maioria das smart TVs e projetores tem suporte; a experiência de assistir séries e filmes melhora muito. Usuários de iPhone também podem usar o AirPlay para espelhar diretamente a tela inteira.
\end{itemize}

\booksection{IV. Leitura de E-books: Também Dá para Ler Bem no Celular}

Tem gente que diz que ler no celular cansa a vista e dispersa a atenção — e isso realmente tem fundamento. Mas as vantagens da leitura no celular também são evidentes — a qualquer hora, em qualquer lugar, aproveitando ao máximo o tempo fragmentado. No metrô, nas filas, antes de dormir, li muitos livros no celular.

\begin{itemize}
  \item \textbf{WeRead (WeChat Reading)}: atualmente, o app de leitura com a melhor experiência na China. O catálogo é muito completo, com livros publicados e web novels. Ele usa o modelo de "tempo de leitura em troca de moedas para livros" — quanto mais você lê, mais livros consegue de graça, o que é muito amigável para estudantes. A função de "Sublinhar" permite ver quais trechos outras pessoas marcaram; às vezes, ver uma boa frase destacada por alguém realmente gera identificação. Mas o "Ranking de Leitura" é um pouco competitivo demais, e pode levar a "ler só para acumular tempo"; sugiro desligar.
  \item \textbf{App Kindle}: se você tem um dispositivo Kindle, o app Kindle no celular sincroniza o progresso de leitura. A qualidade de formatação dos e-books originais do Kindle é muito superior à dos PDFs piratas, mas a loja Kindle China já foi fechada; é preciso migrar para a loja internacional.
  \item \textbf{Apple Books (Livros)}: nativo do iPhone, ideal para ler EPUB e PDF. Interface minimalista, animação de virar página muito agradável. Se você baixou um arquivo de e-book da internet, é só "Compartilhar para o Livros" e ele é importado.
  \item \textbf{Sugestões de configuração de leitura}: independentemente do app, lembre-se de colocar o fundo em tom amarelo quente (modo de proteção ocular), aumentar um pouco a fonte e aumentar o espaçamento entre linhas. Para leitura noturna, ative o "Modo Escuro" — fundo preto com texto branco, muito mais confortável para os olhos. Nunca leia com fundo branco, texto preto e fonte pequena por muito tempo em luz baixa; a fadiga ocular é cumulativa.
\end{itemize}

\booksection{V. Jogos Mobile: Jogar Bem sem Danificar o Celular}

Minha atitude em relação a jogos mobile é: pode jogar, mas não deixe o jogo te controlar. Usar as ferramentas de jogo de forma racional permite jogar com mais conforto e sem atrapalhar as obrigações.

\begin{itemize}
  \item \textbf{Modo Jogo / Espaço de Jogos}: hoje os celulares populares têm essa função. Quando ativado, o celular automaticamente limpa a memória, bloqueia notificações e evita toques acidentais — você está no meio de uma luta em equipe e, de repente, uma chamada aparece; o modo jogo bloqueia isso para você. Alguns celulares também suportam Modo Performance / Modo E-sports, forçando o processador ao máximo, mas o custo é maior aquecimento e consumo de bateria.
  \item \textbf{Como ajustar qualidade gráfica e taxa de quadros}: não coloque tudo no máximo sem pensar. Primeiro garanta a taxa de quadros (fluidez), depois busque qualidade gráfica. 60 fps é muito mais confortável que 30 fps, mas 120 fps pressiona muito o celular; os modelos intermediários não precisam forçar. Se você joga títulos competitivos (como Honor of Kings, PUBG), a prioridade da taxa de quadros é muito maior que a qualidade gráfica — uma jogabilidade fluida é cem vezes mais importante que uma grama bonita.
  \item \textbf{O que fazer se o celular esquentar durante o jogo}: é normal, mas altas temperaturas prolongadas prejudicam a vida útil da bateria. Algumas formas de resfriar: tire a capinha do celular para dissipar calor; reduza o brilho da tela para cerca de 70\% (brilho alto é uma grande fonte de calor); não jogue enquanto carrega (o aquecimento se acumula, a bateria morre rápido). Se o celular estiver quente a ponto de não conseguir segurar, saia imediatamente e deixe-o descansar um pouco.
  \item \textbf{Ferramentas anti-vício}: o "Tempo de Uso" do iOS e o "Bem-Estar Digital" do Android permitem definir limite diário de tempo para apps de jogos. Eu me impus o limite de no máximo uma hora em dias úteis e duas horas nos fins de semana. Quando o tempo acaba, o ícone do app fica cinza e não abre. Esse limite realmente funciona, não é só de fachada.
  \item \textbf{Cloud gaming}: se seu celular não tem desempenho suficiente, mas você quer jogar títulos pesados, experimente uma plataforma de cloud gaming (como Tencent START Cloud Gaming, NetEase Cloud Gaming). O jogo roda num servidor remoto, e a imagem é transmitida como streaming de vídeo para o seu celular. Exige muito da rede, mas quase nada do desempenho do celular. A desvantagem é que a latência é inevitável; é adequado para jogos de estratégia e RPG, não para FPS de tiro.
\end{itemize}

\begin{tcbpractice}{Mãos à Obra: Monte Seu Sistema de Entretenimento Mobile}
\begin{enumerate}
  \item No app de música, baixe seus 3 álbuns ou playlists favoritos e garanta que estejam disponíveis offline.
  \item Vá ao Xiaoyuzhou ou Apple Podcasts, assine 3 canais de seu interesse e baixe um episódio para experimentar.
  \item Instale o VLC Player e tente reproduzir um arquivo de vídeo copiado do computador.
  \item No WeRead, encontre um livro que queira ler, configure fundo em tom quente e modo de proteção ocular, e ajuste a fonte para um tamanho confortável para você.
  \item Se você joga, abra o modo jogo e veja quais configurações podem ser otimizadas (não perturbe, prioridade de taxa de quadros, etc.).
\end{enumerate}

\end{tcbpractice}
